\section*{Summary}

This framework sets out how Visayas State University (\vsu) will use artificial intelligence (AI) in teaching, learning, research, and administration. Its design choices are deliberate:

\begin{itemize}
  \item \textbf{Enable first.} Faculty are encouraged to use AI agents and assistants to design, deliver, and manage their courses, with minimal approval steps and simple disclosure.
  \item \textbf{Proportionate to risk.} Requirements scale with the risk of the use case (a three-tier model for tools; five levels for course-level AI use) rather than blanket bans or uniform gates.
  \item \textbf{Integrity by design.} Academic integrity is protected mainly through clear course-level AI-use declarations, assessment redesign, and process-based evidence, not detector-driven policing.
  \item \textbf{Fair and lawful.} Student data, privacy, equity, and due process are protected in line with the Data Privacy Act and national higher-education policy.
  \item \textbf{Sustainable.} A small governance structure, an annual review cycle, versioned documents, and a feedback loop keep the framework current without heavy bureaucracy.
\end{itemize}

\todoverify{Alignment with \cmo{} and \rmo{} is stated at the level of general intent only. Specific provisions, section numbers, and mandates of those memoranda must be confirmed against the issued texts before this draft advances (see Appendix~\ref{app:verification}).}
